\section{Posentrenamiento y aprendizaje por refuerzo}

\subsection{Elementos centrales del RL}

\textbf{Xie Tianbao (谢天宝) (perspectiva de ingeniería)} enumeró los cuatro elementos centrales del RL:
\begin{enumerate}
  \item \textbf{Environment} — enunciado del problema (conjunto de problemas) + Sandbox (entorno de ejecución). Los requisitos difieren enormemente según el escenario; algunos incluso requieren replicar todo Internet.
  \item \textbf{Foundation Model} — la calidad del Mid-training, el Cold Start SFT y la disposición del propio modelo hacia el RL son todos importantes.
  \item \textbf{Training Infra} — el context de las tareas de RL es notablemente más largo que el de SFT; el diseño de inference engine, context management y agent memory requiere iterar de ida y vuelta.
  \item \textbf{Algoritmo de RL} — el diseño del algoritmo, a su vez, exige coordinación con environment y training infra (por ejemplo, soporte de traceback, optimización de efficiency).
\end{enumerate}
Estos cuatro elementos están altamente acoplados, implican una enorme cantidad de detalles de ingeniería, son difíciles de estudiar por completo en el ámbito académico y también es difícil dominar de una sola vez el know-how.

\begin{importantbox}{Cao Yu (曹玉): el RL es control de lazo cerrado, la definición del problema es más determinante que el algoritmo}
El RL es, en esencia, \textbf{control de lazo cerrado con señal de retroalimentación}; cada eslabón de la cadena de señal es importante. La razón por la que DeepSeek R1 tuvo éxito radica en \textbf{la claridad al elegir y definir el problema} — esto determina el resultado mucho más que el algoritmo concreto.

Los retos reales que enfrenta el RL en la actualidad:
\begin{itemize}
  \item La probabilidad de que el entorno tenga bugs es muy alta
  \item Elegir el problema es muy difícil
  \item Obtener el reward es sparse
  \item Los conflictos entre distintas tareas todavía se resuelven mediante un enfoque human-centric
\end{itemize}
El RL avanza hacia \textbf{la especialización}: quienes trabajan en coding deben convertirse en buenos reward models, y quienes trabajan en verticales (finanzas, derecho, medicina) también deben conocer el dominio.
\end{importantbox}

\subsection{La falta de Scaling Law en el RL}

\textbf{Xue Fuzhao (薛福昭)}: uno de los mayores problemas del RL es \textbf{la falta de un Scaling Law a nivel de código abierto}. Los modelos de código abierto publican versiones ya entrenadas de 7B/14B/70B, no modelos que se ubiquen en la Scaling Curve y permitan predecir por completo el siguiente tier. Esto provoca:
\begin{itemize}
  \item No saber qué ocurrirá con el modelo grande después de entrenar el modelo pequeño
  \item Que la iteración dependa del Run grande final (lento y costoso)
  \item Que el propio loss del RL suba y baje sin control, que el environment sea inestable (actualizaciones de package, problemas de asincronía) y que las conclusiones resulten muy noisy
\end{itemize}
Una vez que se establezca un Scaling Law y un Leaderboard estables, la eficiencia de iteración del RL mejorará enormemente.

\subsection{El problema de la generalización del RL}

\textbf{Cao Yu}: el RL «obtiene exactamente aquello para lo que se entrena, pero eso no significa que mejoren otros aspectos». Actualmente, la magnitud de actualización de parámetros del RL sigue estando en el rango de «pequeña perturbación» (un LoRA rank extremadamente bajo puede igualar al full-parameter), lo que provoca que la capacidad en una sola tarea difícilmente se generalice de forma natural a otros dominios. Pero, en comparación con SFT (offline imitation learning), en teoría el RL debería tener una mejor generalización — la clave es que \textbf{el Scaling del RL todavía está muy lejos de alcanzar la magnitud prometida}.

\textbf{Liu Ziwei (刘子伟)}: no se busca una fuerte generalización hacia tasks completamente unseen; basta con tener cierta \textbf{capacidad de interpolación} dentro de la matriz de tasks conocidos para que ya resulte suficientemente impresionante. Por analogía, la ola de Meta Learning de 2018 tampoco encontró una solución general de generalización. El Scaling de Task y Environment podría ser la clave de la generalización.

\textbf{Xue Fuzhao}: la definición de Generalization no es clara. Si se calcula en FLOPs (cada task muestrea 1024 traces), la generalización del RL hacia un task específico \textbf{debería ser más fuerte} (frente a SFT, que usa un task por cada sequence). El cuello de botella clave es no poder obtener environments a nivel trillion.

\begin{knowledgebox}{Liu Ziwei: distinguir lo verdadero de lo falso}
En los últimos seis meses, el RL se puso muy de moda gracias al reasoning y surgieron muchos trick de variantes. Pero trabajos recientes (como Just RL) han descubierto que muchos de esos trick no son eficaces — \textbf{el enfoque más simple y más limpio, con los hiperparámetros bien ajustados, ya alcanza lo óptimo}. Tras pasar por la bubble, la comunidad empezó a distinguir cuáles avances son reales.
\end{knowledgebox}

\subsection{Resumen de la sección}
El RL pasó de ser un simple adorno a convertirse en la herramienta central para mejorar la capacidad de los modelos, pero enfrenta tres grandes retos: la complejidad del environment, la falta de scaling law y la generalización insuficiente. El éxito de DeepSeek R1 proviene de la claridad en la definición del problema y no de la innovación algorítmica. El RL avanza hacia la especialización.

%% ===== Agentic & Tool Use =====
